\section{Finite quotient witnesses with independent replay}
\label{sec:quotients}

The algebra improvement preserves the exact Khovanov computation. A
complementary strategy is to avoid that computation when a much smaller
obstruction can be found. We implement one such obstruction using $A_5$.
The mathematical method is established; the deliverable here is a bounded
search with explicit certificates, a separate verifier, and a cost
analysis that includes the actual seed-plan construction.

\subsection{Wirtinger relations and soundness}

Divide a knot diagram into Wirtinger arcs, continuing through an
overcrossing and ending at an undercrossing. To each arc assign a group
element. At a crossing, write $u,v$ for the incoming and outgoing
underarc generators and $o$ for the overarc generator. With the sign
$\varepsilon$ determined by the oriented crossing convention, the
relation is
\begin{equation}
 v=o^{\varepsilon}u o^{-\varepsilon}.
 \label{eq:wirtinger}
\end{equation}
The software derives these relations from the validated PD traversal.
A complete satisfying assignment defines a homomorphism from the knot
group into the target group.

\begin{proposition}[Finite-group obstruction]
\label{prop:obstruction}
If a validated one-component knot diagram has a satisfying Wirtinger
assignment with two distinct meridian images in a finite group, the knot
is nontrivial. In fact the image of the representation is nonabelian.
\end{proposition}
\begin{proof}
All oriented meridians of a knot are conjugate in its group. Their images
are therefore conjugate inside the image subgroup. In an abelian subgroup
conjugate elements are equal. Two distinct meridian images thus imply a
nonabelian image. The unknot group is infinite cyclic and has only cyclic
images, so it cannot admit this assignment.
\end{proof}

The search output has only two logical possibilities. A \texttt{KNOTTED}
result includes a witness that has passed independent replay. Every other
result is \texttt{INCONCLUSIVE}, including exhaustive failure for all
chosen targets. The latter proves only the absence of a witness in that
palette. It supplies no evidence that the knot is trivial.

\subsection{Why \texorpdfstring{$A_5$}{A5} and what is stored}

The group $A_5$ is represented by even permutations of five symbols. Its
nonidentity elements have orders $2$, $3$, or $5$. There are fifteen
double transpositions, twenty $3$-cycles, and two conjugacy classes of
twelve $5$-cycles each. Conjugation by an odd permutation is an
automorphism of $A_5$ that exchanges the two $5$-cycle classes. Thus one
chosen $5$-cycle class suffices for an existence search, together with
the other two cycle types. This is a symmetry of the search, not an
assumption that the two classes are equal.

The coefficient field of Khovanov homology and the target group of a
Wirtinger representation play unrelated roles. In particular, an
Alexander-polynomial obstruction can be absent while an $A_5$ witness
exists. Conway and Kinoshita--Terasaka provide concrete fixtures in the
archive. The use of non-affine quandles in that situation has prior
computational and theoretical precedent~\cite{clark,fish}.

A certificate contains explicit permutation images attached to the
diagram, not merely a hash of a search state or a claim that a solver
accepted. The separate checker reconstructs the edge traversal, verifies
the permutation format, checks overarc continuity and every crossing
relation, and checks nonconstancy. It does not import the search's arc
extraction, precomputed conjugation tables, or propagation routines.
The integration evidence includes the precise working PD diagram because
the original pipeline may already have simplified or factored the input.

After a valid diagram has been normalized, replay requires $O(n)$
fixed-size permutation operations and $O(n)$ stored permutation values.
The shipped checker also invokes the existing diagram validation and
normalization, including label sorting. Therefore its entire interface
is not claimed to have literally linear bit time: sorting alone uses
$O(n\log n)$ comparisons on ordinary labels, and label bit lengths must
also be charged.

\section{Certified seeds and exact symmetry reduction}
\label{sec:seeds}

\subsection{A structural propagation plan}

If the overarc and one underarc at a crossing are known,
relation~\eqref{eq:wirtinger} uniquely determines the other underarc.
This observation gives a structural object independent of the target
group.

\begin{definition}[Wirtinger seed plan]
A seed plan consists of $b$ designated arc generators and a sequence of
propagation steps. Each step specifies a crossing whose overarc and one
underarc were designated or produced earlier, and whose other underarc
is new. Every nonseed arc must occur exactly once as a newly produced
arc. A plan is checked by replaying these purely combinatorial
dependencies.
\end{definition}

The existence and optimization of such generating systems belong to the
established Wirtinger-number framework~\cite{blair}. The implementation
constructs a valid plan greedily; it does not claim to find a minimum
seed set or the bridge number.

\begin{theorem}[Complete search from a checked plan]
\label{thm:seedsearch}
Let $C$ be a conjugacy class of size $q$ in a finite group, and let a
validated diagram have a checked $b$-seed plan. Fix the first seed image
to one chosen $c\in C$. Enumerate assignments in $C^{b-1}$ to the
remaining seeds, propagate the plan, and test all crossing relations.
This procedure finds a nonconstant $C$-assignment if and only if one
exists. For a fixed finite target and a supplied plan, its running time
is $O(nq^{b-1})$ and its working storage can be polynomial in $n,b$.
\end{theorem}
\begin{proof}
Every seed assignment extends in at most one way along the plan because
each propagation step is a conjugation or its inverse. All generated
values stay in $C$. Checking every crossing catches relations not used
for propagation and any inconsistent cycles, so an accepted assignment
is a genuine representation. Conversely, every representation restricts
to seed values and is reconstructed from them, by induction over the
plan. Simultaneous conjugation carries its first seed value to $c$, and
preserves the relations and nonconstancy. Thus the normalized enumeration
contains a representative of every possible witness. There are
$q^{b-1}$ tuples; propagation and verification use $O(n)$ group
operations per tuple.
\end{proof}

These operation counts use ordinary arc-index RAM operations and fixed-target
group or table operations. In a strict bit model, charge polynomial factors
in the index bit length as well; they do not change the quasi-polynomial
parameter regimes below.

The greedy plan routine in this archive is conservatively bounded by
$O(n^3)$ for fixed targets and ordinary arc indices. This setup term is
included in the implemented bound. A bound for a supplied plan must not
be advertised as the full interface cost without it.

\subsection{Residual centralizer action}

Fixing one seed color leaves the centralizer $H=Z_{A_5}(c)$ acting by
simultaneous conjugation on the remaining seed tuple. The program
generates a transversal of this action directly. Given a tuple prefix,
it chooses representatives of the current stabilizer's orbits on the next
color, and then replaces the stabilizer by the subgroup fixing that color.
It uses an explicit stack to avoid an artificial recursion-depth limit.

\begin{lemma}[Stabilizer-prefix enumeration]
The recursive orbit construction produces exactly one representative of
each $H$-orbit of complete seed tuples.
\end{lemma}
\begin{proof}
Two tuples in the same $H$-orbit have first entries in the same $H$-orbit,
so choose its unique representative. Once this first entry is fixed, the
elements of $H$ that can identify the remaining suffixes are exactly its
stabilizer. Apply the same argument inductively. This proves existence
and uniqueness at each prefix and hence for complete tuples. The
explicit-stack implementation changes the traversal mechanism but not
the set of prefixes generated.
\end{proof}

For $r=b-1$, Burnside's lemma gives exact counts for the three searches:
\begin{align}
 N_2(b)&=\frac{15^r+3\cdot3^r}{4},\label{eq:N2}\\
 N_3(b)&=\frac{20^r+2\cdot2^r}{3},\label{eq:N3}\\
 N_5(b)&=\frac{12^r+4\cdot2^r}{5}.\label{eq:N5}
\end{align}

\begin{proposition}[Exact $A_5$ seed-orbit counts]
Equations~\eqref{eq:N2}--\eqref{eq:N5} count precisely the complete
tuples generated by the implemented residual inner-centralizer symmetry.
For $b=3$, the respective counts are $63$, $136$, and $32$.
\end{proposition}
\begin{proof}
For an involution $c$, its centralizer is a Klein four-group. The identity
fixes all fifteen involution colors; each of its three nonidentity
elements commutes with exactly three colors in that class. Its action on
$r$-tuples therefore has fixed-point counts $15^r$ and $3^r$.
For a $3$-cycle, the centralizer is cyclic of order three. Each
nonidentity element fixes the two nonidentity powers of that cycle in
the twenty-element class. For a $5$-cycle, the centralizer is cyclic of
order five. Each nonidentity centralizer element fixes exactly the two
powers lying in the chosen twelve-element conjugacy class. Average the
fixed-point counts in each case. Substitute $r=2$ for the numerical
values.
\end{proof}

These formulas refer to this particular symmetry group; they are not a
claim that every possible automorphism symmetry has been exhausted. They
give reproducible maximum trial counts and an independent check of the
tuple generator. For a supplied plan and fixed $A_5$, the seed search
takes $O(nN_j(b))$ time. The complete interface adds its $O(n^3)$ greedy
planning cost. The explicit stack stores tuple prefixes, giving a
conservative auxiliary-space bound $O(n+b^2)$ for fixed $A_5$.

\subsection{Budgets, alternative propagation, and integration}

The seed solver can stop after a prescribed number of complete orbit
assignments or on a cooperative time limit. A separate constraint solver
uses bitset supports, arc consistency, and minimum-domain branching. It
is useful for comparison and can find a witness in a different order,
but its node count and the seed solver's assignment count are distinct
resource measures.

The pipeline extension is optional and runs after the existing invariant
filters, before the expensive homology fallback. Its per-summand time cap
is intersected with the overall remaining time budget. A timeout gives
inconclusive evidence and preserves the original exact fallback. Default
recognition does not enable this search automatically. The measurements
show why: some existing filters already decide the examples in less time.

For general targets, $q^{b-1}$ becomes quasi-polynomial whenever
$b\log(q+1)=O((\log(n+2))^2)$ and target construction, table access, and
plan construction have the corresponding cost bounds. The implemented
$A_5$ case has constant $q$. This is a parameterized coloring bound, not
a universal recognition theorem. In fact, deciding nonconstant coloring
by a fixed nontrivial conjugacy class in a fixed nonabelian finite simple
group is NP-complete~\cite{kuperbergsamperton}. Thus a universal
quasi-polynomial algorithm for that stronger fixed-coloring predicate
would imply $\mathsf{NP}\subseteq\mathsf{QP}$.

\section{An explicit limitation of every small quandle palette}
\label{sec:palette}

The limitation of finite palettes can be proved without relying on any
unproved estimate for the size of a finite quotient. The following
argument also explains why exhausted $A_5$ searches must remain
inconclusive.

\begin{lemma}[Torus power obstruction]
\label{lem:power}
Let $H$ be a finite group of exponent $E$. If $\gcd(p,E)=1$, every
homomorphism
\[
 \langle x,y\mid x^a=y^p\rangle\longrightarrow H
\]
has cyclic image.
\end{lemma}
\begin{proof}
Choose an integer $u$ with $pu\equiv1\pmod E$. For the two generator
images $X,Y$, exponentiation gives
$Y=(Y^p)^u=(X^a)^u=X^{au}$. Both images lie in the cyclic subgroup
generated by $X$.
\end{proof}

A finite quandle $Q$ has bijective right translations
$R_b(a)=a*b$, satisfies $a*a=a$, and satisfies right
self-distributivity. Its inner group $\Inn(Q)$ is the subgroup of
the symmetric group on $Q$ generated by the $R_b$. A quandle coloring
of a knot induces a knot-group representation into $\Inn(Q)$ by sending
meridians to the corresponding right translations. The conjugacy relation
for these translations is precisely the image of a Wirtinger relation.

\begin{lemma}[Cyclic inner image forces a constant knot coloring]
\label{lem:nonfaithful}
If the representation induced by a quandle coloring of a one-component
knot has cyclic image, the coloring is constant. The quandle need not be
faithful or connected.
\end{lemma}
\begin{proof}
The images of knot meridians are conjugate within the image subgroup,
so in a cyclic image all used right translations are the same
permutation $R$. At a crossing with undercolor $a$ and overcolor $b$,
we therefore have $R_b=R_a=R$. Idempotence gives
$R_b(a)=R_a(a)=a$. The inverse translation also fixes $a$. Consequently
the undercolor does not change at a crossing of either sign. The color
also stays fixed along the overarc. Following the single knot component
shows that every color is equal. Equality of right translations alone
would not imply equality of colors in a nonfaithful quandle; the
idempotence and traversal argument is what completes the proof.
\end{proof}

\begin{theorem}[Linear lower bound on detecting palette size]
\label{thm:palette}
For every prime $p>3$, the torus knot $T(3,p)$ is nontrivial and has a
standard $n=2p$-crossing diagram with determinant one. It has no
nonconstant coloring by any finite quandle $Q$ with $|Q|<p$.
Thus a detecting quandle must have at least $n/2$ elements along this
family of inputs.
\end{theorem}
\begin{proof}
The closure of $(\sigma_1\sigma_2)^p$ is a single component because
$\gcd(3,p)=1$, and has $2p$ displayed crossings. The usual decomposition
of its exterior into the two solid-torus pieces gives, by van Kampen,
\[
 G(T(3,p))=\langle x,y\mid x^3=y^p\rangle.
\]
Adding the relation $x^3=y^p=1$ gives the nonabelian quotient
$C_3*C_p$, so the knot group is nonabelian and the knot is nontrivial.

For clarity, its Alexander formula can be read directly from the same
presentation. In abelianization $x\mapsto t^p$, $y\mapsto t^3$. The
two abelianized Fox derivatives are, up to units and sign,
$(t^{3p}-1)/(t^p-1)$ and $(t^{3p}-1)/(t^3-1)$. Their greatest common
divisor is
\[
 \Delta_{3,p}(t)=
 \frac{(t^{3p}-1)(t-1)}{(t^3-1)(t^p-1)},
\]
up to the usual monomial unit, since
$\gcd(t^3-1,t^p-1)=t-1$. Both parameters are odd, so substituting
$t=-1$ gives determinant $|\Delta_{3,p}(-1)|=1$.

Now suppose $|Q|<p$. The order of $\Inn(Q)$ divides $|Q|!$, so its
exponent is coprime to $p$. Lemma~\ref{lem:power} makes the image of
the induced representation cyclic. Lemma~\ref{lem:nonfaithful} then
makes the coloring constant.
\end{proof}

For $A_5$ the exponent is $30$. Every $T(3,p)$ with prime $p\geq7$
therefore has only cyclic $A_5$ images, even when the selected conjugacy
class has more than $p$ elements. The supplied $T(3,7)$ negative control
has a three-seed plan and exhausts exactly $136+63+32=231$ seed-tuple
orbits. Both search implementations return \texttt{INCONCLUSIVE}.

This obstruction does not make those knots hard for the full program.
Their full Alexander polynomial is nontrivial, their diagrams are
homogeneous, and they are three-braid closures. Existing stages can
recognize their nontriviality. The theorem concerns the expressive limit
of a palette alone. It also gives a concrete research question: how large
must a non-affine detecting quandle be for these determinant-one torus
knots? The lower bound does not assert the existence of a detecting
quandle of exactly $p$ elements.
